%% USPSC-Cap3-Desenvolvimento.tex
% ----------------------------------------------------------
% Desenvolvimento
% ----------------------------------------------------------
\chapter{O Ecossistema Trilha das Árvores}
\label{cap:desenvolvimento}

\section{Considerações Iniciais}
\label{sec:dev_consideracoes}

Este capítulo descreve a arquitetura do ecossistema \textit{Trilha das Árvores} e
articula a relação entre as decisões de engenharia de software e a experiência do
usuário delas decorrente. Parte-se da premissa, enunciada na
Seção~\ref{sec:problema}, de que em uma aplicação de uso externo e em movimento a
qualidade percebida da interação não é produto exclusivo do projeto de interface:
ela emerge de toda a pilha tecnológica, da estratégia de persistência de imagens à
política de tratamento de erros de rede.

A descrição a seguir refere-se ao estado do sistema no momento da avaliação, tal
como verificado diretamente no código-fonte dos três repositórios que compõem o
projeto. Onde a implementação atual diverge do que seria a solução arquitetural
desejável, o texto registra a divergência de forma explícita, uma vez que essas
divergências constituem hipóteses verificáveis no estudo empírico e não meros
apontamentos de engenharia.

\section{Visão Geral da Arquitetura}
\label{sec:arquitetura}

O \textit{Trilha das Árvores} adota arquitetura cliente-servidor conteinerizada,
organizada em camadas de responsabilidade delimitadas. O sistema é composto por
dois clientes --- o aplicativo móvel, destinado ao visitante, e o painel
administrativo web, destinado à gestão de conteúdo pela ESALQ --- que se
comunicam por HTTP, com carga útil em JSON, com uma API REST central. Essa API
concentra a lógica de negócio e persiste os dados em um banco relacional.

A adoção do estilo arquitetural REST \cite{fielding2000} confere ao sistema
separação clara entre cliente e servidor e comunicação sem estado, propriedades
que favorecem a evolução independente das interfaces e a previsibilidade do
comportamento observado pelo usuário.

O fluxo de dados percorre quatro estágios. Na \textbf{aquisição}, a camada física
do dispositivo atua como sensor: a câmera lê os QR Codes afixados nas placas das
árvores, o receptor de posicionamento global coleta a localização do usuário e o
magnetômetro fornece a orientação. Na \textbf{apresentação}, o aplicativo móvel
consome esses dados e entrega valor ao visitante, enquanto o painel web executa as
operações administrativas. No \textbf{processamento}, o servidor valida
requisições e aplica as regras de negócio. Na \textbf{persistência}, o banco
relacional garante a durabilidade das informações.

% ============================================================================
% >>> [PREENCHER --- FIGURA DA ARQUITETURA] <<<
% Inserir aqui um diagrama de componentes mostrando: app móvel, painel web,
% API Flask, PostgreSQL, e as integrações (Google Maps no cliente).
% Salvar a imagem em USPSC-img/ e usar o bloco abaixo.
%
% \begin{figure}[htb]
%   \centering
%   \caption{Visão geral da arquitetura do ecossistema Trilha das Árvores}
%   \label{fig:arquitetura}
%   \includegraphics[width=0.95\textwidth]{USPSC-img/arquitetura}
%   \\[4pt]
%   Fonte: Elaborada pelo autor.
% \end{figure}
% ============================================================================

\section{Aplicação Móvel}
\label{sec:app_movel}

\subsection{Organização e fluxo de telas}

O módulo móvel é o principal ponto de contato com o usuário e concentra a maior
parte dos requisitos de experiência. Construído em React Native sobre a
plataforma Expo, organiza-se em telas e componentes reutilizáveis, com navegação
gerenciada por pilha.

O fluxo principal da atividade percorre as seguintes telas: uma tela inicial, que
apresenta o projeto e dá acesso às demais; a listagem de trilhas disponíveis,
obtida da API; uma tela de preparação, que antecede o início e apresenta
distância, número de \textit{checkpoints} e uma lista de verificação prévia; a
tela de atividade, que concentra mapa em tempo real, cronômetro, distância
percorrida e progresso; a tela de leitura de QR Code; e a tela de conclusão. Duas
telas auxiliares completam o conjunto: uma tela informativa sobre o projeto e uma
tela de perfil de usuário.

\subsection{Integração com sensores e recursos de hardware}

A integração com o hardware é realizada por bibliotecas do ecossistema Expo e da
comunidade React Native. O módulo \texttt{expo-location} fornece a geolocalização
contínua que alimenta o posicionamento no mapa; o módulo \texttt{expo-sensors}
provê os dados de magnetômetro que sustentam o componente de bússola; e a
biblioteca \texttt{react-native-vision-camera}, com o módulo de leitura de códigos
habilitado, realiza a decodificação dos QR Codes. A renderização cartográfica
utiliza a biblioteca \texttt{react-native-maps} configurada com o provedor Google
Maps, cuja chave de API é declarada na configuração do projeto.

O cálculo de distância entre pontos apoia-se na fórmula de Haversine, disponível
tanto no cliente quanto no servidor, o que permite exibir ao usuário a distância
acumulada sem depender de requisição de rede a cada atualização de posição.

\subsection{Decisões de interação com impacto sobre a experiência}

Dois aspectos do projeto de interação merecem destaque por seu efeito direto sobre
a experiência.

O primeiro é o \textbf{mecanismo de suspensão e retomada da trilha}. Ao detectar
que o usuário abandona a atividade sem concluí-la, o aplicativo persiste o estado
corrente --- \textit{checkpoint} atual, tempo decorrido, distância acumulada e
lista de árvores --- e oferece, no retorno, a escolha entre retomar do ponto de
interrupção ou reiniciar. O recurso responde a um requisito de robustez diante de
interrupções reais, como uma chamada telefônica, uma perda momentânea de sinal ou
o cansaço, e sustenta a heurística de controle e liberdade do usuário.

O segundo é o \textbf{retorno contínuo de progresso}. Barra de progresso, contagem
de árvores visitadas, cronômetro e distância percorrida permanecem visíveis
durante toda a atividade, e os marcadores no mapa distinguem por cor os
\textit{checkpoints} já visitados, o atual e os subsequentes. Esse conjunto
operacionaliza a heurística de visibilidade do status do sistema ao longo de toda
a execução.

\section{Painel Administrativo}
\label{sec:painel}

O painel administrativo é uma aplicação web de página única construída em React e
empacotada com o conjunto de ferramentas \texttt{react-scripts}, utilizando o
\textit{framework} de estilos Bootstrap para consistência visual e
responsividade. A comunicação com a API é centralizada pela biblioteca Axios, o
que padroniza o tratamento de cabeçalhos de autenticação e de erros.

O painel permite aos gestores da ESALQ manter o conteúdo do parque --- cadastro,
edição e remoção de árvores e de trilhas --- sem necessidade de recompilar ou
redistribuir o aplicativo móvel, desacoplando a operação de conteúdo do ciclo de
publicação do cliente.

Do ponto de vista de interação, destacam-se dois aspectos. O primeiro é o uso da
biblioteca \texttt{@dnd-kit} para ordenação por arrastar e soltar das árvores que
compõem uma trilha, o que torna manipulável de forma direta uma operação que, de
outro modo, exigiria manipulação numérica de índices de ordem.

O segundo é uma decisão de projeto relevante para a integridade dos dados: a
interface opera com o \textbf{identificador institucional da placa} da árvore,
familiar ao administrador que percorre o parque, mas realiza, no momento da
submissão, a tradução para o identificador interno de banco de dados exigido pela
relação de chave estrangeira. Essa tradução ilustra a heurística de
correspondência entre o sistema e o mundo real aplicada à interface
administrativa: expõe-se ao usuário a linguagem que ele conhece, reservando a
complexidade relacional à camada de software.

\section{Servidor de Aplicação}
\label{sec:backend}

\subsection{Organização em camadas}

O servidor é uma API REST desenvolvida em Python com o \textit{microframework}
Flask, estruturada em camadas a fim de reduzir acoplamento e favorecer a
manutenibilidade.

A \textbf{camada de controladores} expõe os \textit{endpoints} e realiza a
validação inicial das requisições, organizada por meio de \textit{Blueprints} que
segmentam o roteamento por domínio funcional: um conjunto de rotas públicas,
consumidas pelo aplicativo móvel, e um conjunto de rotas administrativas,
protegidas por autenticação.

A \textbf{camada de serviços} concentra a lógica de negócio e o encapsulamento de
integrações. Compõem-na o serviço de autenticação, responsável pela emissão e
verificação de tokens; o serviço de banco de dados, que gerencia a sessão; o
serviço de mapa, que implementa o cálculo de distância geodésica pela fórmula de
Haversine; o serviço de configuração e inicialização da aplicação; e um serviço de
armazenamento de objetos, cuja situação é discutida na
Seção~\ref{sec:persistencia}.

A \textbf{camada de modelos} define os esquemas relacionais, mapeados pelo ORM
SQLAlchemy. As entidades principais são \texttt{Tree}, que armazena identificador
institucional, nome comum, nome científico, localização textual, altura, latitude,
longitude e referência ao código de identificação; \texttt{Trail}, que armazena
nome, número de árvores, distância total e imagens; \texttt{Admin}; e a tabela
associativa \texttt{TreeTrail}, que modela o relacionamento entre árvores e
trilhas e carrega dois atributos próprios da associação --- a ordem da árvore no
percurso e a distância em relação à árvore anterior.

\subsection{Confiabilidade da conexão com o banco}

Uma decisão de implementação de interesse particular para a confiabilidade
percebida é a ativação da opção \texttt{pool\_pre\_ping} nas opções de motor do
SQLAlchemy, realizada na inicialização da aplicação. A opção determina que cada
conexão obtida do \textit{pool} seja testada antes do uso, o que previne erros de
conexão expirada após períodos de inatividade do banco. Tais erros, se propagados
ao cliente, manifestar-se-iam como mensagens obscuras e interrupção de fluxo
durante a trilha. Trata-se de exemplo concreto de configuração de infraestrutura,
invisível ao usuário, que sustenta diretamente a heurística de prevenção de erros.

\subsection{Documentação da interface de programação}

A documentação da API é servida interativamente por meio de Swagger, com a
interface disponível em rota dedicada. Registra-se, contudo, que a especificação
publicada encontra-se \textbf{desatualizada} em relação às rotas efetivamente
implementadas, e que seus textos descritivos conservam conteúdo de exemplo do
andaime inicial do projeto. A correção da especificação figura entre as
recomendações consolidadas no Capítulo~\ref{cap:conclusao}.

\section{Persistência de Dados e de Artefatos Binários}
\label{sec:persistencia}

Os dados estruturados são persistidos em um sistema gerenciador de banco de dados
relacional PostgreSQL, escolhido por sua robustez, conformidade transacional e
maturidade no ecossistema de código aberto.

O tratamento dos artefatos binários merece registro detalhado, por constituir
divergência entre a arquitetura projetada e a implementada. O projeto \textbf{prevê}
um serviço de armazenamento de objetos compatível com a API do Amazon S3,
implementado sobre MinIO e encapsulado em um serviço dedicado, cuja classe está
presente no código-fonte e cuja dependência consta das exigências do projeto.
Entretanto, na versão atualmente em operação, esse serviço \textbf{não é invocado
por nenhum controlador}: as imagens de capa e de mapa das trilhas são armazenadas
como colunas binárias na própria tabela de trilhas e servidas por
\textit{endpoints} do processo Flask.

Essa divergência tem duas consequências relevantes para a experiência do usuário.
A primeira é que o servimento de imagens ocorre pelo processo de aplicação, que
não é otimizado para entregar arquivos estáticos volumosos de forma concorrente;
sob carga, isso eleva a latência de carregamento e prolonga telas de espera no
cliente. A segunda é que o armazenamento de dados binários em colunas do banco
relacional onera consultas e rotinas de cópia de segurança, o que compromete a
escalabilidade da solução à medida que o acervo cresce.

Registra-se ainda que a configuração do serviço de armazenamento contém endereço
de rede fixado em código, acompanhado de anotação de pendência --- indício de que
a integração permaneceu incompleta. A conclusão dessa integração figura entre as
recomendações de evolução prioritária.

\section{Serviços Cartográficos}
\label{sec:cartografia}

A base cartográfica e a plotagem dos marcadores apoiam-se no provedor Google Maps,
consumido \textbf{diretamente pelo cliente móvel} por meio da biblioteca
\texttt{react-native-maps}, com chave de API declarada na configuração do
aplicativo. O servidor não intermedia requisições cartográficas: seu serviço de
mapa limita-se ao cálculo geodésico de distâncias.

Essa configuração tem três implicações. Primeiro, a exibição do mapa depende de
conectividade no dispositivo do usuário; em áreas de mata densa do campus, a
latência ou a indisponibilidade da rede degrada a orientação, o que reforça a
importância dos mecanismos de resiliência no cliente e de estados de carregamento
e erro bem comunicados. Segundo, não há previsão de operação com base cartográfica
armazenada localmente, de modo que a ausência de conectividade compromete a
funcionalidade central da atividade. Terceiro, a chave de API encontra-se exposta
em arquivo versionado do repositório, o que constitui vulnerabilidade de
configuração e recomenda migração para variável de ambiente com restrição de uso
por aplicativo.

\section{Segurança e Autenticação}
\label{sec:seguranca}

O acesso às operações administrativas é protegido por autenticação baseada em
\textit{JSON Web Tokens}, implementada com a biblioteca PyJWT e centralizada no
serviço de autenticação. O modelo sem estado do JWT alinha-se ao estilo REST e
simplifica a verificação de autorização a cada requisição.

No plano da experiência, o tratamento diferenciado das respostas de erro ---
distinguindo, por exemplo, ``não autorizado'' de ``erro nos dados enviados'' ---
permite ao cliente apresentar mensagens específicas e acionáveis, atendendo à
heurística de ajuda ao usuário no reconhecimento e na recuperação de erros.

Registram-se, contudo, três limitações de segurança identificadas na
implementação atual. A primeira é que a chave de assinatura dos tokens está
fixada em código-fonte versionado, o que compromete a confidencialidade da chave e
inviabiliza sua rotação. A segunda é que os tokens emitidos não contêm campo de
expiração, de modo que um token, uma vez emitido, permanece válido
indefinidamente. A terceira é que a comunicação com o servidor institucional
ocorre sobre HTTP em texto claro, o que exigiu configuração explícita de
permissão de tráfego não cifrado no cliente Android. Embora funcional no contexto
controlado da rede da universidade, a ausência de TLS constitui limitação de
segurança e de confiança percebida, e sua correção figura como evolução
prioritária.

\section{Implantação e Infraestrutura}
\label{sec:implantacao}

O servidor de aplicação e o banco de dados são conteinerizados e orquestrados com
Docker Compose, em execução sobre servidor Linux na infraestrutura de rede da
universidade. A composição define dois serviços --- o banco PostgreSQL, com volume
persistente e verificação de saúde, e a API, que depende da disponibilidade
efetiva do banco para iniciar --- e expõe a API em porta dedicada, com política de
reinício automático.

A conteinerização garante paridade entre os ambientes de desenvolvimento e de
produção, isolando dependências do sistema hospedeiro e reduzindo a classe de
defeitos associada a divergências de ambiente. Do ponto de vista da avaliação, essa
previsibilidade contribui para que o software submetido ao estudo com usuários
corresponda ao software efetivamente entregue.

\section{Decisões Técnicas como Suporte à Experiência do Usuário}
\label{sec:mapeamento}

A discussão precedente sustenta a tese de que, em aplicações móveis de uso em
campo, a experiência do usuário é propriedade emergente de toda a pilha
tecnológica, e não apenas da camada de apresentação. O Quadro~\ref{quad:mapeamento}
consolida essa correlação, associando cada decisão de engenharia ao efeito
esperado sobre a percepção do usuário e à heurística correspondente, conforme a
numeração adotada no Quadro~\ref{quad:heuristicas}.

\begin{quadro}[htb]
\centering
\caption{Mapeamento entre decisões técnicas, efeito sobre a experiência e heurística associada}
\label{quad:mapeamento}
\footnotesize
\begin{tabular}{p{4.0cm} p{4.6cm} p{2.6cm}}
\toprule
\textbf{Decisão técnica} & \textbf{Efeito esperado sobre a experiência} &
\textbf{Heurística} \\
\midrule
React Native com acesso nativo a sensores & Baixa latência de posicionamento,
bússola e câmera; imersão preservada & H2 \\
\addlinespace
Componentes de interface reutilizáveis & Consistência visual e comportamental
entre telas & H4 \\
\addlinespace
Suspensão e retomada da trilha & Robustez a interrupções; senso de controle sobre
a atividade & H3 \\
\addlinespace
Retorno contínuo de progresso, tempo e distância & Usuário permanentemente
informado do estado da atividade & H1 \\
\addlinespace
Marcadores diferenciados por cor no mapa & Distinção imediata entre visitado,
atual e subsequente & H1, H6 \\
\addlinespace
Tradução de identificador de placa para identificador interno & Administrador
opera na linguagem que conhece & H2 \\
\addlinespace
\texttt{pool\_pre\_ping} no pool de conexões & Prevenção de erros de conexão
expirada & H5 \\
\addlinespace
Tratamento diferenciado de códigos de erro HTTP & Mensagens de erro específicas e
acionáveis & H9 \\
\midrule
\multicolumn{3}{l}{\textit{Limitações identificadas}} \\
\midrule
Servimento de imagens pelo processo Flask & Risco de latência sob carga; telas de
espera prolongadas & H1 (negativa) \\
\addlinespace
Dependência de conectividade para a base cartográfica & Risco de indisponibilidade
do mapa em áreas de mata & H5, H9 (negativa) \\
\addlinespace
Comunicação sem TLS & Redução da confiança e da segurança percebidas & --- \\
\addlinespace
Token sem expiração e chave em código & Risco de segurança não percebido pelo
usuário final & --- \\
\addlinespace
Tela de perfil com dados não integrados & Informação irreal; quebra de confiança
no sistema & H1 (negativa) \\
\addlinespace
Especificação da API desatualizada & Não afeta o usuário final; onera a evolução &
--- \\
\bottomrule
\end{tabular}
\\[4pt]
\begin{flushleft}
Fonte: Elaborado pelo autor.
\end{flushleft}
\end{quadro}

As linhas assinaladas como limitações não constituem apenas apontamentos de
engenharia: são \textbf{hipóteses verificáveis} no estudo com usuários. Espera-se
que se manifestem como atritos observáveis sob o protocolo \textit{Think-Aloud} e
que encontrem correspondência nos escores de Facilidade de Uso Percebida do
instrumento TAM. É esse encadeamento --- da decisão técnica ao atrito observado, e
do atrito observado ao escore atitudinal --- que a matriz de rastreabilidade
descrita no capítulo seguinte se destina a explicitar.

\section{Considerações Finais}
\label{sec:dev_finais}

Este capítulo descreveu a arquitetura do ecossistema \textit{Trilha das Árvores}
em seus três módulos, detalhou as decisões de engenharia com potencial impacto
sobre a experiência do usuário e registrou, de forma explícita, as divergências
entre a arquitetura projetada e a efetivamente implementada, bem como as
limitações de segurança e de desempenho identificadas na versão avaliada. O
mapeamento consolidado no Quadro~\ref{quad:mapeamento} cumpre o primeiro objetivo
específico enunciado na Seção~\ref{sec:objetivos} e fornece o conjunto de
hipóteses que a avaliação empírica, descrita no capítulo seguinte, se propõe a
verificar.
